%%%%%
%
\section{Maximum-Echantillonnage type rafale d'un param\`etre {\tt float}}

\itemstyle{Op\'erateur~:} {\sf MaxEch\_raf}
\index{MaxEch\_raf}

%%%%%OBJET%%%%
\itemstyle{Objet~:}
Calcul des valeurs maxi par \'echantillon d'un
param\`etre \`a une fr\'equence donn\'ee, le temps
en sortie correspond \`a la fin de l\'echantillon
(et non au temps du maxi comme dans MaxEch)

%%%%%ARGUMENTS%%%%
\bigskip
%%%%%ENTREES%%%%
\begin{tabarg}
\multicolumn{3}{l}{\titretab {Entr\'ees}}\\
\hline
$X$ & Vecteur, {\tt float} & Param\`etre d'entr\'ee\\
$Time-in$ & Vecteur, {\tt double} & Vecteur temps correspondant au vecteur d'entr\'ee\\
$R_f$ & Coeff., {\tt double} & Rapport des fr\'equences ($ F_{sortie}\over
F_{entree} $) \\
$Decal$ & Coeff., {\tt int} & D\'ecalage (\textit{nombre de points})\\
\end{tabarg}

\medskip
%%%%%SORTIES%%%%
\begin{tabarg}
\multicolumn{3}{l}{\titretab {Sorties}}\\
\hline
$Time-out$ & Vecteur, {\tt double} & Vecteur temps correspondant a la fin de l\'echantillon\\
$X_m$ & Vecteur, {\tt float} & R\'esultat (m\^eme unit\'e que $X$)\\
\end{tabarg}


%%%%%FORMULE%%%%
\itemstyle{Formule}

%%Max
$$
Recherche de la valeur maxi dans chaque echantillon
\qquad i=0,Nbre\_Echant(X). R_f
$$


%%%%COMMENTAIRE%%%%
\itemstyle{Commentaire~:}

$R_f$ doit \^etre compris dans l'intervalle $]0, 1]$,
et \^etre l'inverse d'un entier. \newline
Le nombre d'\'echantillons du vecteur de sortie est le produit du
nombre d'\'echantillons du vecteur d'entr\'ee par $R_f$, arrondi
\`a l'entier sup\'erieur.

%%

%%%%%AUTEUR%%%%
\itemstyle{Auteur~:}
CNRM/GMEI/TRAMM

%%%%%REFERENCE%%%%
\itemstyle{R\'ef\'erences~:}
